\documentclass[10pt,a4paper]{amsart}
\usepackage[margin=19mm]{geometry}
\usepackage{amsmath,amssymb,mathtools}
\usepackage[scr=rsfs,cal=euler]{mathalfa}
\usepackage{xfrac}
\usepackage[unicode,hidelinks,bookmarks=false]{hyperref}
\newcommand{\ie}{i.e.\ }
\newcommand{\cf}{cf.\ }
\newcommand{\resp}{respectively\ }
\newcommand{\Subsection}{\subsection}
\providecommand{\nobreakdash}{\nobreak\hbox{-}}
\setlength{\emergencystretch}{2em}
\multlinegap0pt
\usepackage{amssymb}
\usepackage{mathtools}
\usepackage{tikz}
\usepackage[shortlabels]{enumitem}
\newtheorem{theorem}{Theorem}
\newtheorem{claim}{Claim}
\newcommand{\cA}{\mathcal{A}}
\newcommand{\cC}{\mathcal{C}}
\newcommand{\del}{\!\setminus\!}
\newcommand{\ep}{\varepsilon}
\title{Matroids from gain graphs over quotient groups}
\author{Zach Walsh}
\date{Source: arXiv:2602.23066v1 (CC BY 4.0)}
\begin{document}
\maketitle

\noindent\emph{Source and licence.} Matroids from gain graphs over quotient groups. Version arXiv:2602.23066v1 (CC BY 4.0), distributed by arXiv under the Creative Commons Attribution 4.0 International licence, http://creativecommons.org/licenses/by/4.0/, as recorded in the arXiv licence metadata for this exact version and shown on the abstract page https://arxiv.org/abs/2602.23066v1. Full licence notice: \texttt{licenses/arxiv-2602.23066-CC-BY-4.0.txt}.\par\smallskip

\noindent\textit{Editorial scope.} Counted unit: the article's theorem precise (source0.tex lines 936--939). The article's complete original proof of it is delivered with it (source0.tex lines 940--1058, one \texttt{proof} environment of 119 lines). No mathematics is written, repaired or re-derived: the statement and the proof are the article's own lines, in the article's own notation, and no retained line is substituted or re-flowed.  No further source-local context is needed: every cross-reference of every retained line resolves inside the two delivered spans.  Nothing was omitted from any retained span, so the package carries no omission record.  No retained line contains a citation, so the wrapper carries no bibliography and no bibliography entry.  No retained line contains a figure environment, an image-inclusion command or drawing code, so no image is delivered and the package declares no asset dependency.  Editorial formatting only, in addition: an AMS \texttt{amsart} preamble that restates the article's own declarations and macros used by the retained lines, namely the environments \texttt{theorem}, \texttt{claim} on separate counters, together with the standard packages those lines need.  The retained environments print in this document as Theorem 1, Claim 1 in order of appearance, whereas the article prints the same items with the numbering of its own full document.  The complete native source of the article as distributed in the arXiv e-print for this exact version is bundled unchanged as \texttt{sources/195367/source0.tex}.
\begin{theorem} \label{thm: main precise}
Let $\Gamma$ be a group, let $\Gamma_1$ be a normal subgroup of $\Gamma$, let $\{\Gamma_1\} \cup \cA$ be a partition of $\Gamma$ so that each $A \in \cA$ is malnormal and has all conjugates in $\cA$, and let $(G, \psi)$ be a $\Gamma$-gain graph.
Then $\cC(\Gamma, \Gamma_1, \cA, G, \psi)$ is a linear class of circuits of the underlying $(\Gamma/\Gamma_1)$-frame matroid $F(G, \psi/\Gamma_1)$.
\end{theorem}

\begin{proof}
We will write $N$ for the frame matroid $F(G, \psi/\Gamma_1)$, $\cC$ for $\cC(\Gamma, \Gamma_1, \cA, G, \psi)$, and $\ep$ for the identity element of $\Gamma$.
Let $C_1, C_2 \in \cC$ with $|C_1 \cup C_2| = r_N(C_1 \cup C_2) + 2$, and let $C$ be a circuit of $N$ contained in $C_1 \cup C_2$ with $C \notin \{C_1, C_2\}$.
We will show that $C \in \cC$.
We first prove two important properties of the graph $G[C_1 \cup C_2]$.
The \emph{cyclomatic number} of a graph is the minimum number of edges whose deletion leaves a forest.

\begin{claim} \label{claim: cyclomatic number 3}
The graph $G[C_1 \cup C_2]$ is connected and has cyclomatic number at most $3$.
\end{claim}
\begin{proof}
Note that $G[C_1]$ and $G[C_2]$ are both connected graphs.
Then $G[C_1 \cup C_2]$ is connected or else $C \in \{C_1, C_2\}$.
Since $|C_1 \cup C_2| = r_N(C_1 \cup C_2) + 2$, we can delete two edges from $C_1 \cup C_2$ to obtain an independent set $I$ of the frame matroid $N$.
This implies that each component of $G[I]$ has at most one cycle.
If $G[I]$ is connected, then the claim follows.
If $G[I]$ is disconnected, then it has at most three connected components because it is obtained by deleting two edges from a connected graph.
If it has three components, then neither deletion decreased the cyclomatic number and the claim follows.
If it has two components, then one of the two deletions did not decrease the cyclomatic number and the claim follows.
\end{proof}


If $G[C_1 \cup C_2]$ has cyclomatic number $2$ then $C_1$ and $C_2$ are both $\psi$-balanced cycles.
If they share no edges, then $C \in \{C_1, C_2\}$.
If they share an edge then $G[C_1 \cup C_2]$ is a theta graph, and it follows from the theta property for $\psi$-balanced cycles that $C$ is a $\psi$-balanced cycle and is therefore in $\cC$.
So for the remainder of the proof we may assume that $G[C_1 \cup C_2]$ has cyclomatic number $3$.
The next claim deals with $\psi$-balanced cycles.





\begin{claim} \label{claim: no balanced cycles}
If $G[C_1 \cup C_2]$ contains a $\psi$-balanced cycle, then $C \in \cC$.
\end{claim}
\begin{proof}
Suppose $G[C_1 \cup C_2]$ contains a $\psi$-balanced cycle $D$.
Let $T$ be a spanning tree of $G[C_1 \cup C_2]$ that contains all but one edge of $D$.
By switching, we may assume that $T$ is $\psi$-normalized.
Since $D$ is $\psi$-balanced, it follows that $D$ is $\psi$-normalized.
Since $G[C_1 \cup C_2]$ has cyclomatic number $3$, there are two edges $e$ and $f$ in $(C_1 \cup C_2) - (T \cup D)$.
We will first show that there is some $A \in \cA$ that contains $\psi(e)$ and $\psi(f)$.
If $C_1$ is a handcuff or a theta graph then $e,f \in C_1$, and since $C_1 \in \cC$ there is some $A \in \cA$ that contains $\psi(e)$ and $\psi(f)$.
The same reasoning applies if $C_2$ is a handcuff or a theta graph, so we may assume that $C_1$ and $C_2$ are both $\psi$-balanced cycles.
By choosing $D = C_1$ we may assume that $C_1$ is $\psi$-normalized.
Then $e$ and $f$ are both in $C_2$, and since $C_2$ is $\psi$-balanced and every edge other than $e$ and $f$ is $\psi$-normalized it follows that $\psi(e) = \psi(f) = \{\alpha, \alpha^{-1}\}$ for some $\alpha \in \Gamma - \{\ep\}$.
If $\alpha \in \Gamma_1$ then $C_1 \cup C_2$ is $(\psi/\Gamma_1)$-normalized and so $r_N(C_1 \cup C_2) = |V(G[C_1 \cup C_2])| - 1 = |C_1 \cup C_2| - 2 - 1$ because $G[C_1 \cup C_2]$ has cyclomatic number $3$.
But then $|C_1 \cup C_2| - r_N(C_1 \cup C_2) = 3$, a contradiction.
So $\alpha \notin \Gamma_1$, so there is some $A \in \cA$ with $\alpha \in A$, and therefore $A$ contains $\psi(e)$ and $\psi(f)$.



If $C$ is a handcuff or a theta graph, then $e,f \in C$ and it follows that $C \in \cC$.
So we may assume that $C$ is a $(\psi/\Gamma_1)$-balanced cycle.
We must show that $C$ is $\psi$-balanced.
If $C$ contains neither $e$ nor $f$, then $C$ is $\psi$-normalized.
If $C$ contains exactly one of $e$ and $f$, then $C$ is not $(\psi/\Gamma_1)$-balanced because $\alpha \notin \Gamma_1$, a contradiction.
So $C$ contains $e$ and $f$.
Let $W$ be a simple closed walk around $C$.
Then $\psi(W)$ is the product of an element in $\psi(e)$ and an element in $\psi(f)$, so $\psi(W) \in A$.
And $\psi(W) \in \Gamma_1$ because $C$ is $(\psi/\Gamma_1)$-balanced.
So $\psi(W) = \ep$ because $A \cap \Gamma_1 = \{\ep\}$, and therefore $C$ is $\psi$-balanced and in $\cC$, as desired.
\end{proof}



By Claim \ref{claim: no balanced cycles} we may assume that $G[C_1 \cup C_2]$ contains no $\psi$-balanced cycles. 
This implies that $C_1$ and $C_2$ are not $(\psi/\Gamma_1)$-balanced; if $C_i$ were $(\psi/\Gamma_1)$-balanced, it would be $\psi$-balanced because it is in $\cC$.
Note that $C$ may be $(\psi/\Gamma_1)$-balanced.
We consider one more special case.

\begin{claim} \label{claim: only one cycle}
If $G[C_1 \cap C_2]$ has at most one cycle, then $C \in \cC$.
\end{claim}
\begin{proof}
Suppose that $G[C_1 \cap C_2]$ has at most one cycle.
Then there is an edge $e_1 \in C_1 - C_2$ that is in a cycle of $G[C_1]$ and an edge $e_2 \in C_2 - C_1$ that is in a cycle of $G[C_2]$.
Moreover, $G[C_1] \del e_1$ and $G[C_2] \del e_2$ are both connected, and since they share at least one vertex it follows that $G[C_1 \cup C_2] \del \{e_1, e_2\}$ is connected.
Since $G[C_1 \cup C_2]$ has cyclomatic number $3$ we see that $G[C_1 \cup C_2] \del \{e_1, e_2\}$ has exactly one cycle.
Let $e_3$ be an edge of this cycle, and note that $G[C_1 \cup C_2] \del \{e_1, e_2, e_3\}$ is a spanning tree of $G[C_1 \cup C_2]$.
By switching we may assume that it is $\psi$-normalized.
Note that for $i = 1,2,3$, $\psi(e_i) \ne \{\ep\}$ because $G[C_1 \cup C_2]$ contains no $\psi$-balanced cycles.
Since $C_1$ contains $e_1$ and $e_3$ but not $e_2$ and $C_1$ is a handcuff or theta graph, there is some $A_1 \in \cA$ that contains $\psi(e_1)$ and $\psi(e_3)$ since $C_1 \in \cC$.
Similarly, there is some $A_2 \in \cA$ that contains $\psi(e_2)$ and $\psi(e_3)$.
Then $A_1 = A_2$ because both contain $\psi(e_3)$.
Let $T$ be a spanning tree of $G[C]$.
By switching within $A_1$ we can normalize $T$ so that the edge(s) in $C - T$ have gain values in $A_1$.
If $C$ is a theta graph or a handcuff this implies that $C \in \cC$, so we may assume that $C$ is a $(\psi/\Gamma_1)$-balanced cycle.
Let $W$ be a simple closed walk around $C$.
Clearly $\psi(W) \in A_1$.
Also, $\psi(W) \in \Gamma_1$ since $C$ is $(\psi/\Gamma_1)$-balanced.
Therefore $\psi(W) = \ep$ since $A_1 \cap \Gamma_1 = \{\ep\}$, so $C \in \cC$.
\end{proof}

By Claim \ref{claim: only one cycle} we may assume that $G[C_1 \cap C_2]$ contains at least two cycles.
And by Claim \ref{claim: no balanced cycles} we may assume that $C_1$ and $C_2$ are not cycles.
Since $C_1$ has cyclomatic number two it follows that $G[C_1 \cap C_2]$ has cyclomatic number two.
Therefore for $i = 1,2$, $C_i$ is a non-minimal graph with cyclomatic number two, so $C_1$ and $C_2$ are loose handcuffs that contain the same pair of cycles, say $D$ and $D'$.
For $i = 1,2$, let $P_i$ be the path in $G[C_i]$ between $V(D)$ and $V(D')$.
Let $v$ be the end of $P_2$ in $V(D)$.
Let $e$ and $f$ be edges of $D$ and $P_2$, respectively, that are incident with $v$.
Let $e'$ be any edge of $D'$.
By switching we may assume that the spanning tree $(C_1 \cup C_2) - \{e, e', f\}$ of $G[C_1 \cup C_2]$ is $\psi$-normalized.
Since $C_1$ contains $e$ and $e'$ but not $f$ and $C_1 \in \cC$, there is some $A_1 \in \cA$ that contains $\psi(e)$ and $\psi(e')$.
Let $\{w,v\}$ and $\{u,v\}$ be the sets of ends of $e$ and $f$, respectively.
Let $\psi'$ be the $\Gamma$-gain function obtained from $\psi$ by switching by $\psi(f, u, v)^{-1}$ at each vertex in $V(D)$.
Then $\psi'(f) = \ep$ and $\psi'(e') = \psi(e')$, while $\psi'(e, v, w) = \psi(f, u, v) \circ \psi(e, v, w) \circ \psi(f, u, v)^{-1}$ and $\psi'(g) = \{\ep\}$ for each edge $g$ in $C_2 - \{e, e'\}$.
Since $C_2 \in \cC$, there is some $A_2 \in \cA$ that contains $\psi'(e')$ and $\psi'(e, v, w)$.
Since $\psi'(e') = \psi(e')$ we see that $A_2 = A_1$.
Since $A_1$ contains $\psi(e, v, w)$ and $\psi(f, u, v) \circ \psi(e, v, w) \circ \psi(f, u, v)^{-1}$ and $A_1$ is malnormal it follows that $\psi(f, u, v) \in A_1$.
So $A_1$ contains $\psi(e) \cup \psi(e') \cup \psi(f)$, and therefore $A_1$ contains $\psi(C_1 \cup C_2)$.
Let $T$ be a spanning tree of $G[C]$.
By starting with $\psi$ and switching within $A_1$ we can $\psi$-normalize $T$ so that the edge(s) in $T - C$ have gain values in $A_1$.
If $C$ is a theta graph or a handcuff this implies that $C \in \cC$, so we may assume that $C$ is a $(\psi/\Gamma_1)$-balanced cycle.
Let $W$ be a simple closed walk around $C$.
Clearly $\psi(W) \in A_1$.
Also, $\psi(W) \in \Gamma_1$ since $C$ is $(\psi/\Gamma_1)$-balanced.
Therefore $\psi(W) = \ep$ since $A_1 \cap \Gamma_1 = \{\ep\}$, so $C \in \cC$.
\end{proof}

\end{document}
